\documentclass[11pt]{article}
\usepackage[a4paper,margin=1in]{geometry}
\usepackage[T1]{fontenc}
\usepackage[utf8]{inputenc}
\usepackage{lmodern}
\usepackage{amsmath,amssymb,amsthm}
\usepackage{microtype}
\usepackage{hyperref}
\title{Heinrich Weber\\Lehrbuch der Algebra\\English Translation Draft\\Group Characters and Reciprocal Groups}
\author{}
\date{}
\begin{document}
\maketitle

\section*{\S 11. Group characters}

A principal problem in the theory of Abelian groups is to determine all divisors of an Abelian group. The solution of this problem is substantially facilitated by introducing the concept of \emph{group characters}, which is also important in many other respects.

Let \(S\) be an Abelian group of order \(n\), whose elements we denote by \(A\), and let
\[
A_1,A_2,\dots,A_\nu
\tag{1}
\]
be the elements of a basis of \(S\), of orders
\[
a_1,a_2,\dots,a_\nu.
\]
Then every element \(A\) is represented, and represented uniquely, in the form
\[
A = A_1^{\alpha_1}A_2^{\alpha_2}\cdots A_\nu^{\alpha_\nu},
\tag{2}
\]
where
\[
0\le \alpha_1<a_1,\qquad 0\le \alpha_2<a_2,\qquad \dots,\qquad 0\le \alpha_\nu<a_\nu.
\tag{3}
\]
The exponents \(\alpha_1,\alpha_2,\dots,\alpha_\nu\), or any integers congruent to them modulo \(a_1,a_2,\dots,a_\nu\), are called the \emph{indices} of the element \(A\). Then one has:

\medskip
\noindent\textbf{1.} The indices of a product \(AA'\) are obtained by adding the corresponding indices of the two factors.

\medskip
If we assign to each element \(A\) some numerical value, we may call this assignment a function of \(A\). Such a function \(\chi(A)\) will be called a \emph{group character} if for every two elements \(A,A'\in S\) it satisfies
\[
\chi(AA')=\chi(A)\chi(A').
\tag{4}
\]
Hence it also satisfies the more general multiplicative law
\[
\chi(A'A''A'''\cdots)=\chi(A')\chi(A'')\chi(A''')\cdots.
\]

Two such functions \(\chi\) and \(\chi_1\) are regarded as distinct if there exists at least one element \(A\) for which \(\chi(A)\neq \chi_1(A)\). We now determine these characters more closely.

Setting \(A'=1\) in \((4)\), one obtains
\[
\chi(A)=\chi(A)\chi(1).
\tag{5}
\]
Excluding once and for all the degenerate case in which all \(\chi(A)=0\), one gets:
\[
\chi(1)=1.
\tag{6}
\]

\medskip
\noindent\textbf{2.} Every character takes the value \(1\) on the identity element.

\medskip
Setting \(A'=A\) in \((4)\), we obtain
\[
\chi(A^2)=[\chi(A)]^2,
\]
and by complete induction, for every exponent \(h\),
\[
\chi(A^h)=[\chi(A)]^h.
\tag{7}
\]
If \(a\) denotes the order of the element \(A\), then with \(h=a\) and using \((6)\) one gets
\[
[\chi(A)]^a=1.
\tag{8}
\]

\medskip
\noindent\textbf{3.} The values of a character \(\chi(A)\) are roots of unity whose orders divide the order of \(A\).

\medskip
We can now determine all characters easily. Set
\[
\chi(A_1)=\omega_1,\qquad \chi(A_2)=\omega_2,\qquad \dots,\qquad \chi(A_\nu)=\omega_\nu.
\tag{9}
\]
Then
\[
\omega_1^{a_1}=1,\qquad \omega_2^{a_2}=1,\qquad \dots,\qquad \omega_\nu^{a_\nu}=1,
\tag{10}
\]
and from \((2)\) and \((4)\),
\[
\chi(A)=\omega_1^{\alpha_1}\omega_2^{\alpha_2}\cdots \omega_\nu^{\alpha_\nu}.
\tag{11}
\]
Conversely, whenever \(\omega_1,\omega_2,\dots,\omega_\nu\) is any solution of \((10)\), formula \((11)\) defines a function of \(A\) that satisfies \((4)\), and hence is a character of the group.

Each of the equations \((10)\) has \(a_1,a_2,\dots,a_\nu\) roots, and if all root choices are combined with one another, one obtains
\[
a_1a_2\cdots a_\nu=n
\]
possible systems. These lead to \(n\) distinct characters \(\chi(A)\). Indeed, if another choice
\[
\omega'_1,\omega'_2,\dots,\omega'_\nu
\]
yields the same function for all \(A\), then taking \(\alpha_1=1\), \(\alpha_2=\cdots=\alpha_\nu=0\) shows \(\omega_1=\omega'_1\), and similarly for the others.

Thus:

\medskip
\noindent\textbf{4.} There exist exactly \(n\) distinct characters of an Abelian group of order \(n\), all given by formula \((11)\).

\medskip
Let \(\varepsilon_1,\varepsilon_2,\dots,\varepsilon_\nu\) be primitive roots of the equations \((10)\). Then we may write
\[
\omega_1=\varepsilon_1^{\beta_1},\qquad \omega_2=\varepsilon_2^{\beta_2},\qquad \dots,\qquad \omega_\nu=\varepsilon_\nu^{\beta_\nu},
\tag{12}
\]
where \(\beta_1,\dots,\beta_\nu\) are taken from complete residue systems modulo \(a_1,\dots,a_\nu\). Then all \(n\) characters are given by
\[
\chi(A)=\varepsilon_1^{\alpha_1\beta_1}\varepsilon_2^{\alpha_2\beta_2}\cdots\varepsilon_\nu^{\alpha_\nu\beta_\nu}.
\tag{13}
\]
Among these occurs the character obtained by setting all \(\beta_i=0\); it has value \(1\) on every element \(A\), and we call it the \emph{unit character}.

The \(n\) characters of \(S\) can themselves be combined into an Abelian group, indeed into one isomorphic to \(S\).

For by \((13)\), each character is determined by a system of integers
\[
\beta_1,\beta_2,\dots,\beta_\nu
\]
taken modulo \(a_1,\dots,a_\nu\), and this same system is the index system of a definite element \(B\in S\), namely
\[
B=A_1^{\beta_1}A_2^{\beta_2}\cdots A_\nu^{\beta_\nu}.
\tag{14}
\]
Thus every element \(B\in S\) corresponds to a definite character, which we denote by \(\chi_B(A)\), or simply \(\chi_B\) when the variable \(A\) is omitted.

If \(B'\) is a second element of \(S\), with indices \(\beta'_1,\beta'_2,\dots,\beta'_\nu\), then its corresponding character \(\chi_{B'}\) is defined likewise. If now \(\chi_B\chi_{B'}\) denotes the character given on each \(A\) by
\[
\chi_B\chi_{B'}(A)
=
\varepsilon_1^{\alpha_1(\beta_1+\beta'_1)}
\varepsilon_2^{\alpha_2(\beta_2+\beta'_2)}
\cdots
\varepsilon_\nu^{\alpha_\nu(\beta_\nu+\beta'_\nu)},
\tag{15}
\]
then the characters are thereby combined into a group isomorphic to \(S\). The identity element of this character group is the unit character.

Accordingly, the character group can itself be represented by a basis, obtained by setting
\[
\chi_1(A)=\varepsilon_1^{\alpha_1},\qquad
\chi_2(A)=\varepsilon_2^{\alpha_2},\qquad
\dots,\qquad
\chi_\nu(A)=\varepsilon_\nu^{\alpha_\nu}.
\tag{16}
\]
Then every character has the form
\[
\chi_B = \chi_1^{\beta_1}\chi_2^{\beta_2}\cdots \chi_\nu^{\beta_\nu},
\tag{17}
\]
and if \(B,B'\in S\), then
\[
\chi_B\chi_{B'}=\chi_{BB'}.
\tag{18}
\]
If \(\chi_1,\chi_2\) are any two characters and \(\chi_1\chi_2\) their composite, then for every \(A\),
\[
\chi_1(A)\chi_2(A) = (\chi_1\chi_2)(A).
\tag{19}
\]

We may summarize this as:

\medskip
\noindent\textbf{5.} The characters of a group \(S\) can be combined into a group isomorphic to \(S\).

\medskip
One has also, immediately from definition \((13)\),
\[
\chi_B(A)=\chi_A(B).
\tag{20}
\]

Finally there is the following theorem.

\medskip
\noindent\textbf{6.} If \(A\) runs through all elements of the group \(S\), then for a fixed character \(\chi\),
\[
\sum_A \chi(A)=
\begin{cases}
n,& \text{if }\chi\text{ is the unit character},\\[4pt]
0,& \text{otherwise}.
\end{cases}
\tag{21}
\]
Likewise, if the element \(A\) is held fixed and \(\chi\) runs through all characters, then
\[
\sum_\chi \chi(A)=
\begin{cases}
n,& \text{if }A=1,\\[4pt]
0,& \text{otherwise}.
\end{cases}
\tag{22}
\]

When \(\chi\) is the unit character, or \(A\) is the identity, these formulas are evident, since each of the \(n\) summands is then equal to \(1\). If \(\chi\) is not the unit character, there exists \(B\in S\) such that \(\chi(B)\neq 1\). Writing
\[
\sum_A \chi(A)=\sigma,
\]
and multiplying by \(\chi(B)\), one obtains
\[
\sum_A \chi(AB)=\sigma\chi(B).
\]
But as \(A\) runs through the whole group, so does \(AB\), hence the left side is also \(\sigma\). Therefore
\[
\sigma[1-\chi(B)]=0,
\]
and since \(\chi(B)\neq 1\), we get \(\sigma=0\), proving \((21)\). Formula \((22)\) follows similarly, and also immediately from \((20)\).

\section*{\S 12. Divisors of an Abelian group. Reciprocal groups}

A divisor \(T\) of an Abelian group \(S\) is, as already remarked above, always a normal subgroup. Hence by \S 4 there exists a complementary group
\[
S/T,
\]
whose elements are the cosets of \(T\), namely
\[
T,\quad TA',\quad TA'',\quad \dots,
\tag{1}
\]
where \(A',A'',\dots\) are certain elements of \(S\). If \(t\) is the order of \(T\), then the number of such cosets, that is, the order of the group \(S/T\), is equal to the index of \(T\) in \(S\), namely
\[
\frac{n}{t}.
\]

This group \(S/T\) is itself Abelian, because
\[
TA'\cdot TA'' = TA'A'',
\qquad
TA''\cdot TA' = TA''A',
\]
and these are equal.

Let us now determine the characters of the group \(S/T\). Denote its elements by
\[
T_0,T_1,T_2,\dots,T_j,
\tag{2}
\]
and let one of its characters be written \(\xi(T_i)\).

From this function \(\xi(T)\) we can derive a function \(\xi(A)\) on the elements of \(S\) by setting
\[
\xi(A_i)=\xi(T_i)
\tag{3}
\]
whenever \(A_i\) is any element belonging to the coset \(T_i\). For the elements of the group \(T\) itself one then has \(\xi(A)=1\).

But this function \(\xi(A)\) is among the characters of \(S\). For if \(A_i\in T_i\) and \(A_k\in T_k\), then \(A_iA_k\in T_iT_k\), and hence
\[
\xi(A_i)\xi(A_k)=\xi(T_i)\xi(T_k)=\xi(T_iT_k)=\xi(A_iA_k).
\tag{4}
\]
This is precisely the defining property of a character of \(S\).

Conversely, suppose one of the characters \(\chi=\xi\) of \(S\) takes the same value on every element of \(T\). Since the identity belongs to \(T\), this common value can only be \(1\). If now \(A_i\) is any element of the coset \(T_i\), and \(A\) runs through the whole group \(T\), then \(AA_i\) runs through the coset \(T_i\). Since \(\xi(A)=1\), one has
\[
\xi(AA_i)=\xi(A_i).
\tag{5}
\]
That is, \(\xi(A)\) has one and the same value on all elements of a fixed coset \(T_i\), and may therefore be regarded as a function of \(T_i\), denoted \(\xi(T_i)\).

Since whenever \(A_i\in T_i\) and \(A_k\in T_k\), the product \(A_iA_k\in T_iT_k\), it follows that
\[
\xi(T_i)\xi(T_k)=\xi(T_iT_k).
\tag{6}
\]
So \(\xi(T_i)\) is indeed a character of the quotient group \(S/T\).

Therefore there are exactly \(j\) such characters \(\xi(A)\), namely those that take the value \(1\) on every element of \(T\). These \(j\) characters form, under composition of characters, a group; for if \(\xi_1(A)=1\) and \(\xi_2(A)=1\) on \(T\), then also \(\xi_1\xi_2(A)=1\). Since these functions \(\xi\) can equally well be viewed as the characters of the quotient group \(S/T\), it follows from \S 11, theorem 5, that the group of the \(\xi\) is isomorphic to \(S/T\).

Hence:

\medskip
\noindent\textbf{7.} If an Abelian group \(S\) has a divisor \(T\) of order \(t\) and index \(j\), then among the characters of \(S\) there are exactly \(j\), and no more, that take the value \(1\) on all elements of \(T\); and for every element of \(S\) not contained in \(T\), at least one of these characters takes a value different from \(1\). These characters form a group isomorphic to \(S/T\).

\medskip
Since every character \(\chi_B\) corresponds to a definite element \(B\in S\), if \(\chi_B\) runs through the group of characters belonging to \(T\), then \(B\) runs through a group of order \(j\), isomorphic to the group of the \(\chi_B\), and therefore also to \(S/T\). We denote this group by \(U\) and call it the \emph{group reciprocal to \(T\)}.

Again, one should note that in general the notion of reciprocal group depends on the chosen basis and the chosen roots of unity \(\varepsilon\), so it does not belong intrinsically to the group \(S\) and the subgroup \(T\) alone.

The reciprocal group of \(T\) is characterized by the following theorem.

\medskip
\noindent\textbf{8.} Let \(A\) run through the elements of a divisor \(T\subset S\). Then the elements \(B\in S\) satisfying
\[
\chi_B(A)=1
\tag{7}
\]
for every \(A\in T\) are precisely the elements of the group \(U\) reciprocal to \(T\), whose order is equal to the index of \(T\).

\medskip
By \S 11, \((20)\), the group \(T\) is in turn the reciprocal group to \(U\); the relation between the two groups is therefore mutual.

One also has the following theorem.

\medskip
\noindent\textbf{9.} If \(T'\) is a divisor of \(T\), then conversely the reciprocal group \(U\) of \(T\) is a divisor of the group \(U'\) reciprocal to \(T'\).

\medskip
For every element \(A'\in T'\) also lies in \(T\). Hence, if \(B\) runs through \(U\), then \(\chi_B(A')=1\), and therefore \(B\) must occur in the group \(U'\).

If one selects from the totality of characters \(\chi\) an arbitrary finite number
\[
\xi_1,\xi_2,\dots,\xi_k,
\]
whether or not they form a group, then all elements \(A\) satisfying the \(k\) conditions
\[
\xi_1(A)=1,\qquad
\xi_2(A)=1,\qquad
\dots,\qquad
\xi_k(A)=1
\tag{8}
\]
form a group, because from \(\xi_i(A)=1\) and \(\xi_i(A')=1\) it follows that \(\xi_i(AA')=1\).

If the characters \(\xi_1,\dots,\xi_k\) do not already form a group, then there follow from equations \((8)\) additional equations
\[
\xi_{k+1}(A)=1,\quad \dots
\]
so that \(\xi_1,\dots,\xi_k\) may be completed to a group.

\medskip
\noindent\textbf{10.} It follows from theorem 7 that in this way one can obtain all possible divisors of \(S\).

\medskip

\end{document}
